\begin{table}[h]
  \begin{center}
    \topcaption{Sculpting ratio $R(m_{a})$ of Eq.~(\ref{eq:sculpt_R}), evaluated on the
    inclusive simulated background, for the low-mass and high-mass BDT models. $R=1$
    corresponds to no sculpting of the $m_{\ell\ell\gamma\gamma}$ spectrum in the
    Higgs-peak window; $R>1$ ($R<1$) indicates an enhancement (depletion) of the
    peak window relative to the sidebands. The largest value,
    $R(m_{a}=30~\GeV)=1.30$, is not a concern: $R$ is evaluated on the simulated
    background, whereas the final background is the data-driven sideband fit, and a
    single fixed-window ratio cannot distinguish a broad, smoothly-varying tilt of the
    mass shape (absorbed by the analytic background function) from a localized fake
    peak. This is confirmed by the background-only pseudo-data closure test
    (Fig.~\ref{fig:sculpt_closure}), whose observed limit at $m_{a}=30~\GeV$ lies within
    the $\pm1\sigma$ expected band with a best-fit signal strength compatible with zero,
    showing that the $R>1$ there does not bias the result.}
    \label{tab:sculpt_R}
    \small
    \begin{tabular}{|c|c||c|c|} \hline
      \multicolumn{2}{|c||}{Low-mass model} & \multicolumn{2}{c|}{High-mass model} \\ \hline
      $m_{a}$ [\GeV] & $R(m_{a})$ & $m_{a}$ [\GeV] & $R(m_{a})$ \\ \hline
      1 & 0.87 & 4  & 0.44 \\
      2 & 0.96 & 5  & 0.34 \\
      3 & 1.21 & 6  & 0.53 \\
        &      & 7  & 0.82 \\
        &      & 8  & 1.01 \\
        &      & 9  & 0.99 \\
        &      & 10 & 0.77 \\
        &      & 15 & 1.11 \\
        &      & 20 & 0.98 \\
        &      & 25 & 1.12 \\
        &      & 30 & 1.30 \\ \hline
    \end{tabular}
  \end{center}
\end{table}
